% Two insertions. The first is a new subsection for Appendix B (the parameterisation control).
% The second replaces the Related Work treatment of source conditioning in Section 2.

% =====================================================================
% INSERT AS A NEW SUBSECTION OF APPENDIX B, after the capacity sweep.
% =====================================================================

\subsection{The gap is not an artefact of the velocity parameterisation}

$\mathcal M_2$ regresses the velocity. Given $(\bm z_t,\bm z_0)$ the paired target is exactly
recoverable, and Proposition~3 shows that recovery has a $1/t$ Jacobian. A natural objection is
that the degradation is therefore an artefact of asking the network to represent a quantity whose
implicit inversion is ill-conditioned at small $t$, and that predicting the endpoint directly
would remove it.

We test this. In the endpoint parameterisation the network predicts $\hat{\bm z}_1$ as a residual
from the current state, $\hat{\bm z}_1=\mathrm{normalize}(\bm z_t+f_\theta(\bm z_t,t,\bm c,\bm a))$,
is trained against $\bm z_1$ directly, and induces the sampling velocity
$\log_{\bm z}(\hat{\bm z}_1)/(1-t)$, so that the geodesic reaches the predicted endpoint at
$t=1$. The residual form is used because a zero-initialised output layer makes
$\mathrm{normalize}(f_\theta)$ degenerate at initialisation; the remaining time is floored at half
an integration step so that the terminal Heun evaluation stays finite. Everything else --- the
architecture, the optimizer, the schedule, the seeds and the evaluation --- is unchanged, and the
evaluation set is seeded so that the two parameterisations are scored on identical conditions.
Five seeds per cell.

\begin{center}
\begin{tabular}{@{}llrrrrrr@{}}
\toprule
parameterisation & $\eta$ & $\mathcal M_1$ & $\mathcal M_2$ & gap & gap/se &
$\mathcal M_1$ pair & $\mathcal M_2$ pair\\
\midrule
velocity & $10^{-4}$ & $0.3170$ & $0.2614$ & $+0.0556$ & $48.8$ & $0.202$ & $0.437$\\
endpoint & $10^{-4}$ & $0.2906$ & $0.2374$ & $+0.0532$ & $80.0$ & $0.264$ & $0.500$\\
\addlinespace
velocity & $10^{-3}$ & $0.3149$ & $0.1932$ & $+0.1217$ & $127.7$ & $0.213$ & $0.773$\\
endpoint & $10^{-3}$ & $0.2844$ & $0.1995$ & $+0.0850$ & $21.5$ & $0.213$ & $0.639$\\
\bottomrule
\end{tabular}
\end{center}

The gap survives. At $\eta=10^{-4}$ it retains $96\%$ of its size under endpoint prediction, at
$80$ standard errors; at $\eta=10^{-3}$ it retains $70\%$; averaged over the two rates, $78\%$.
The conditional-diversity collapse survives as well, with $\mathcal M_2$ roughly twice
$\mathcal M_1$'s pairwise cosine under either parameterisation. And the inversion that motivates
the analysis is unchanged: under endpoint prediction $\mathcal M_2$ still attains the
\emph{lower} training loss at both learning rates ($4.09\times10^{-4}$ against
$5.78\times10^{-4}$, and $3.69\times10^{-4}$ against $5.61\times10^{-4}$), while producing the
worse rollout.

We report one qualification. Retention is near-complete at the smaller learning rate and
incomplete at the larger one, so while the gap itself is structural, roughly $30\%$ of the
\emph{additional} gap induced by a larger step size is specific to the velocity
parameterisation. The claim supported by this experiment is that source conditioning degrades the
rollout under either parameterisation, not that the parameterisation is irrelevant to the
magnitude.

% =====================================================================
% REPLACES / EXTENDS THE "RECOVERING COUPLING INFORMATION" PARAGRAPH IN SECTION 2.
% =====================================================================

\paragraph{Where the source enters, in practice.}
Two successful methods make opposite choices about the informative source, and the distinction
this paper isolates is exactly the one that separates them.

Posterior-Mean Rectified Flow \citep{ohayon2024pmrf} predicts the posterior mean and then
transports it to the target distribution with a rectified flow. Its source is as informative as a
source can be, and yet the transport improves the result --- because the objective is the
distribution rather than the pointwise reconstruction, and because the flow's velocity field is
\emph{not} conditioned on the source: the source only sets the initial condition. In our
terminology PMRF is $\mathcal M_1$. It also adds Gaussian noise to the posterior mean before
training, and states the reason directly: the source and target lie on manifolds of different
dimension, and the noise alleviates the singularity of learning a deterministic map between them.
That is the same obstruction we formalise in Corollary~\ref{cor:degen} and measure in
Appendix~\ref{sec:srdim}, and the same remedy we study in Appendix~\ref{sec:remedies}, applied
there to the initialisation rather than to the auxiliary input.

Augmented Bridge Matching \citep{debortoli2023abm} makes the other choice: it supplies the initial
sample to the velocity field, explicitly trading the Markov property for the ability to recover
the training coupling. Its objective is coupling recovery for mixture-of-translation tasks, not
conditional sample quality, and on that objective it succeeds. In our terminology it is
$\mathcal M_2$.

Our contribution is not that either is wrong. It is the controlled measurement of what that
architectural choice costs when the objective is the conditional distribution: at matched
capacity, training and initialisation, supplying the source to the field lowers the regression
objective and worsens the rollout, in all twelve learning-rate by width cells, under both the
velocity and the endpoint parameterisation. Read this way, the result explains why PMRF is built
the way it is, and delimits the regime in which the augmentation of \citet{debortoli2023abm} is
the right construction.

% Suggested bibliography entries -----------------------------------------
% @article{ohayon2024pmrf,
%   title  = {Posterior-Mean Rectified Flow: Towards Minimum MSE Photo-Realistic
%             Image Restoration},
%   author = {Ohayon, Guy and Michaeli, Tomer and Elad, Michael},
%   journal= {arXiv preprint arXiv:2410.00418}, year = {2024}}
% @inproceedings{blau2018perception,
%   title  = {The Perception-Distortion Tradeoff},
%   author = {Blau, Yochai and Michaeli, Tomer},
%   booktitle = {CVPR}, year = {2018}}
% Also add: Freirich et al. (theoretical DP curve / optimal estimator),
%           Kim et al. 2026 (Better Source Better Flow, arXiv:2602.05951),
%           Albergo et al. 2023 (stochastic interpolants),
%           Delbracio & Milanfar (InDI), Tsimpos et al., Zhang et al. (BGM),
%           Pooladian et al. (Multisample Flow Matching).
